\chapter*{Conclusion générale}
\addcontentsline{toc}{chapter}{Conclusion générale}

\section*{Bilan}
\addcontentsline{toc}{section}{Bilan}

Ce projet de fin d'études a abouti à la réalisation de \textbf{PharmacieConnect}, une plateforme web et
mobile répondant à un besoin concret : faciliter l'accès aux informations sur les pharmacies et les
pharmacies de garde, tout en proposant un espace de gestion pour les administrateurs et une assistance
conversationnelle pour les premiers secours. La solution s'appuie sur une API backend, un tableau de bord
administrateur, une application mobile et un assistant de premiers secours basé sur la récupération
d'informations hybride.

\section*{Résultats obtenus}
\addcontentsline{toc}{section}{Résultats obtenus}

Les objectifs principaux ont été atteints. L'utilisateur mobile peut rechercher des pharmacies, utiliser la
géolocalisation, consulter les pharmacies de garde par date et accéder au catalogue de médicaments. Il
peut également consulter l'assistant de premiers secours pour obtenir des réponses fiables et sûres à des
questions d'urgence médicale. L'administrateur peut gérer les pharmacies, les gardes, les médicaments,
les utilisateurs autorisés, les imports CSV, les indicateurs et les journaux d'audit. Le backend assure la
centralisation des données, l'authentification, le contrôle des rôles, la validation des API et la traçabilité
des actions sensibles. L'assistant offre une récupération précise d'informations avec un Recall@5 de 94.38\%,
une fusion hybride BM25/sémantique et des mécanismes de confiance et de sécurité pour éviter les
hallucinations et ajouter les mises en garde appropriées pour les urgences.

Sur le plan technique, ce travail a mobilisé plusieurs compétences : conception d'API REST,
développement backend avec FastAPI, modélisation SQLAlchemy, développement web avec React et Vite,
développement mobile avec Expo et React Native, architecture de récupération d'informations hybride,
entraînement et optimisation de pipelines NLP, authentification JWT, validation de fichiers CSV,
journalisation, indicateurs et tests automatisés. Il a également montré l'intérêt d'une séparation claire
entre les routes, les services, les modèles, les interfaces clientes, la base de données et les services
spécialisés comme l'assistant.

\section*{Limites}
\addcontentsline{toc}{section}{Limites}

Certaines limites restent à prendre en compte. Le catalogue des médicaments ne représente pas encore le
stock réel de chaque pharmacie. Les données de garde dépendent toujours de la qualité des fichiers
importés. L'assistant, bien que performant sur la récupération, pourrait bénéficier d'une amélioration de la
sélection de spans (F1 = 0.2990 contre une cible de 0.35), possibilité de réentraînement avec des
données d'entraînement enrichies pour affiner davantage l'extraction. Les tests bout en bout et les tests
de charge doivent encore être renforcés. L'application mobile devra aussi être testée sur davantage
d'appareils et dans différents contextes de connexion. L'intégration de l'assistant au sein de l'application
mobile nécessite des tests supplémentaires pour assurer la cohérence de l'expérience utilisateur.

\section*{Perspectives}
\addcontentsline{toc}{section}{Perspectives}

Les perspectives d'amélioration sont les suivantes :

\begin{itemize}
    \item ajouter la disponibilité des médicaments par pharmacie ;
    \item permettre aux pharmacies de mettre à jour directement certaines informations sous contrôle ;
    \item envoyer des notifications lorsque les gardes changent ;
    \item améliorer la recherche par synonymes, erreurs de frappe ou noms commerciaux proches ;
    \item améliorer l'extraction de spans de réponse de l'assistant grâce à un reclassement plus fin ou un rétiquetage des données ;
    \item étendre la base de connaissance de l'assistant à d'autres sujets médicaux pertinents pour les pharmacies ;
    \item ajouter des capacités multilingues robustes à l'assistant (tests sur l'arabe et l'anglais) ;
    \item intégrer des tests bout en bout couvrant l'application mobile, le tableau de bord, l'API et l'assistant ;
    \item renforcer le déploiement par une configuration de production, une supervision, des sauvegardes et un monitoring de l'assistant ;
    \item enrichir les indicateurs pour mieux suivre les recherches, la couverture territoriale et l'utilisation de l'assistant ;
    \item construire une interface d'administration pour gérer les données d'entraînement et les performances de l'assistant.
\end{itemize}

Ainsi, PharmacieConnect constitue une base fonctionnelle pour une solution complète de gestion
pharmaceutique avec assistance intelligente. Le projet peut encore évoluer, mais il montre déjà qu'une
architecture bien organisée, associée à des techniques modernes comme la récupération hybride
d'informations, peut répondre à la fois aux besoins du citoyen, aux besoins de l'administration et aux
besoins d'assistance médicale d'urgence.
